\documentclass[11pt]{article}
\usepackage[margin=1in]{geometry}
\usepackage{amsmath,amssymb,amsthm}
\usepackage{booktabs}
\usepackage{graphicx}
\usepackage{natbib}
\usepackage{hyperref}
\hypersetup{colorlinks=true,linkcolor=blue,citecolor=blue,urlcolor=blue}

\newtheorem{proposition}{Proposition}
\newtheorem{assumption}{Assumption}
\newtheorem{definition}{Definition}

\title{Fragmented Housing Regulation as a Fiscal-Federalism Breakdown\\
\large The Distortion from Local Free-Riding, and the Leakage of State Preemption}
\author{Daniel Posthumus}
\date{June 2026}

\begin{document}
\maketitle

\begin{abstract}
\noindent Local housing regulation is provided by jurisdictions far finer than the labor
market and housing submarkets they sit in. Each jurisdiction enjoys the local amenity of low
density but does not internalize the regional cost of the housing supply it suppresses --- the
classic condition (\citealp{oates1972fiscal}; \citealp{besley2003centralized}) under which
decentralized provision of a local public good with interjurisdictional spillovers ceases to
be efficient. This note frames fragmented housing regulation as exactly that breakdown and
builds a small equilibrium model of the \emph{market for regulation} aimed at two questions a
purely descriptive account cannot answer. \textbf{First, how large is the distortion} from
decentralized, non-internalized regulation, relative to a regional (internalized) benchmark?
\textbf{Second, what is the leakage of state preemption} --- when the state forces open one
regulatory margin (the by-right zoning envelope), to what extent do jurisdictions reassert
restriction through the discretionary margin it leaves untouched? The model's central
empirical object is the cross-jurisdiction \emph{reaction function}: in an efficient,
no-spillover world jurisdictions would not respond to one another's regulation, so a nonzero
strategic-interaction slope is the direct signature of the spillover that breaks the
Decentralization Theorem. Every primitive maps to an item-level dataset built from
planning-commission minutes paired with the by-right envelope, across Bay Area jurisdictions.
\end{abstract}

\tableofcontents

\section{The question, and why it is a federalism question}

A large literature establishes that local land-use regulation suppresses housing supply in
high-demand metros at large aggregate cost \citep{glaeser2005,hsieh2019housing,glaeser2009causes}.
Recent work measures the \emph{by-right} regulatory environment at scale, including with LLMs
\citep{bartik2024costs,gyourko2021local}. What that work measures is a level (or a set of
levels) of stringency; what it largely takes as given is \emph{why} regulation is set where it
is, and what would happen under a different \emph{assignment of regulatory authority}. The
``homevoter'' and ``jurisdictional fragmentation'' hypotheses supply the political ingredients
\citep{fischel2004economic,fischel2001homevoter,elmendorf2025homevoter,ortalo2014political},
but not a tractable equilibrium object that prices the externality and supports counterfactuals.

\paragraph{The federalism frame.} The Decentralization Theorem \citep{oates1972fiscal} holds
that local provision of a local public good is efficient \emph{absent interjurisdictional
spillovers and scale economies}; \citet{besley2003centralized} recast the centralization
trade-off as a political-economy problem, where provision is chosen by elected local
representatives (decentralization) or a legislature (centralization), and the comparison turns
on spillovers and taste heterogeneity operating through political channels. Housing regulation
is precisely the case the theorem's condition excludes. A jurisdiction that restricts
development consumes a \emph{local} amenity --- low density, preserved ``character,'' a
scarcity rent on incumbent housing --- while the cost, suppressed housing supply, is borne
\emph{regionally}: by would-be entrants priced out of the metro, by the labor market that loses
agglomeration when workers cannot locate near high-productivity jobs, and by every other
jurisdiction in the same housing submarket whose prices the restriction props up. The benefit
is internalized; the cost spills across boundaries onto people who do not vote in the
restricting jurisdiction. Decentralized provision is therefore inefficient, and the inefficiency
is not incidental --- it is the structural reason fragmented housing regulation over-restricts.

\paragraph{What this buys, and the two questions.} Framing the distortion as a federalism
breakdown gives it a \emph{named efficiency benchmark} (the regional, internalized optimum) and
turns two otherwise-vague goals into sharp questions.
\begin{enumerate}
\item \textbf{The magnitude of the distortion.} How much regulation, and how much suppressed
housing supply, is attributable to fragmentation per se --- the gap between regulation chosen by
atomistic localities each ignoring the spillover, and regulation chosen at the level that
internalizes it (the submarket or the region)? This is the Oates wedge, made quantitative.
\item \textbf{The leakage of preemption.} State preemption is \emph{partial centralization}: it
forces open one regulatory margin and leaves another decentralized. If a jurisdiction regulates
on two substitutable margins --- a \emph{by-right envelope} (zoning, the margin preemption
targets) and a \emph{discretionary friction} (planning-commission review, which it retains) ---
then preemption that loosens the envelope can be partly or wholly offset by tightening
discretion. The leakage is the fraction of the intended loosening that localities claw back
through the untouched margin. A descriptive measure of either margin alone cannot see it.
\end{enumerate}
\textbf{What is new here} is the combination: a two-margin structure that makes preemption
leakage a coherent object, embedded in a fragmented system whose strategic interaction is the
empirical signature of the very spillover that makes decentralization inefficient.

\section{Environment}\label{sec:env}

One metro --- the Bay Area --- comprises $J$ jurisdictions $j=1,\dots,J$, grouped into $S\ll J$
housing submarkets $s(j)$, all drawing on \emph{one} regional labor market. A worker employed
anywhere in the metro earns the regional wage $w_m=w(N)$, $w'>0$, $N\equiv\sum_j N_j$, capturing
agglomeration rents \citep{hsieh2019housing}. This is the coupling that makes regulation's cost
regional: a household's income depends on the metro, not on the jurisdiction it sleeps in, so the
benefit of access is regional while the disamenity of the housing that provides it is local.

\paragraph{Preferences and prices.} A household of type $\theta$ in $j$ has
$U=c+\alpha\ln h-\gamma g(D_j)$, $g'>0$, where $-\gamma g(D_j)$ is the \emph{local} disamenity of
density (so low density is a normal good and high-$\theta$ residents are the natural demanders of
restriction), and income $y(\theta)=w_m\theta$. Within submarket $s$, clearing gives the price
index
\begin{equation}
p_{s} = \frac{\alpha N_{s}}{H_{s}}, \qquad
N_s=\!\!\sum_{j:s(j)=s}\!\!N_j,\ \ H_s=\!\!\sum_{j:s(j)=s}\!\!H_j,
\label{eq:price}
\end{equation}
so housing built in any jurisdiction of $s$ lowers the price faced by incumbents in \emph{every}
jurisdiction of $s$; density $D_j=N_j/\bar L_j$ is local. (We treat $h$ as a normalized unit of
housing services so that $H_s$ counts dwellings consistent with the data; the divisible-services
reading and the indivisible-dwelling/spatial-indifference reading coincide under that
normalization, and nothing below depends on the distinction.) Equations
\eqref{eq:price} contain the externality: the \emph{cost} of approval (density) is local, the
\emph{benefit} (lower submarket price, more regional employment hence higher $w_m$) is shared.

\section{Regulation on two margins}\label{sec:tech}

The empirical premise, visible in the minutes: among projects that reach a vote, approval (or
its discretionary-review equivalent) is nearly universal. Restriction does not operate through
denial at the hearing; it operates before and around it, on two margins.

\paragraph{The by-right envelope $\bar z_j$.} The zoning map assigns each parcel a use and
height/bulk district, defining the scale/use a project may obtain \emph{ministerially} --- with
certainty, no hearing. Projects within the envelope never appear on the commission calendar;
this is the margin LLM-based zoning measurement captures \citep{bartik2024costs}, and the margin
state preemption targets.

\paragraph{The discretionary friction $\tau_j$.} A project exceeding the envelope (a use not
principally permitted, a scale above the limit, a special-district trigger) enters discretionary
review and is built only after a \emph{discretionary tax} $\tau_j$ whose components the minutes
record: process exposure (the share of activity forced off the by-right path, via
\texttt{request\_type}), mobilized opposition ($g(D_j)$ made operative, via stance-marked speaker
counts), conditioning (\texttt{action}$=$approved-with-conditions), delay (continuances, via a
case-level join), and rare denial. This is the margin LLM zoning measurement \emph{cannot} see
and the margin preemption leaves untouched.

\begin{definition}[Effective regulation]
Effective stringency is $r_j=g(\bar z_j,\tau_j)$, increasing in the tightness of the envelope
(smaller $\bar z_j$ pushes more activity into discretion) and in the discretionary tax. A
jurisdiction restricts by lowering $\bar z_j$ \emph{or} raising $\tau_j$; the two are
substitutes in producing restriction but face different costs, and --- crucially --- state
preemption acts on $\bar z_j$ but not $\tau_j$.
\end{definition}

\section{The homevoter's choice and the fragmentation externality}\label{sec:pol}

Each jurisdiction's regulation is chosen by a vote-maximizing authority answering to incumbent
homeowners (share $\lambda_j$, weight $\phi_j$), following the homevoter logic
\citep{fischel2004economic}. Mapping the two margins to two bodies resolves the unitary-authority
shortcut: the council sets the slow map $\bar z_j$, the planning commission sets the fast
discretionary tax $\tau_j$. An incumbent's payoff, taking neighbors' policies as given, is
\begin{equation}
W^{\text{own}}_j = \underbrace{p_{s(j)}(\bar z_{-j},\tau_{-j};\bar z_j,\tau_j)}_{\text{submarket asset value}}
 - \underbrace{\gamma\,g\!\big(D_j(\bar z_j,\tau_j)\big)}_{\text{local disamenity}},
\label{eq:owner}
\end{equation}
where $-j$ denotes the profile of all other jurisdictions (only co-submarket neighbors,
$s(k)=s(j)$, enter with nonzero weight). The authority maximizes
$\Omega_j=\phi_j\lambda_j W^{\text{own}}_j+(1-\phi_j\lambda_j)W^{\text{ent}}_j$ over
$(\bar z_j,\tau_j)$.

\paragraph{The free-riding wedge.} The effect of $j$'s own approvals on the price its homeowners
care about is attenuated by submarket size: $\partial p_{s(j)}/\partial(\text{supply}_j)
=-(p_s/H_s)$. A jurisdiction small within its submarket bears the full \emph{local} disamenity of
any approval but captures little of the price relief, most of which leaks to co-submarket
neighbors. So it under-supplies and free-rides on their construction --- the microfoundation of
the fragmentation externality.

\begin{proposition}[Fragmentation over-restricts; the distortion is the Oates wedge]\label{prop:wedge}
Holding $(\phi_j\lambda_j,\gamma,\bar L_j)$ fixed, each jurisdiction's privately optimal
restriction is increasing in the number of jurisdictions sharing its submarket: as $j$ becomes
atomistic the price-relief motive vanishes, the local-disamenity motive dominates, and both
$\bar z_j$ tightens and $\tau_j$ rises. The \textbf{distortion} is the gap between the
decentralized Nash regulation $\{(\bar z_j^\ast,\tau_j^\ast)\}$ and the regulation chosen by an
authority that internalizes the submarket (or region) --- the regional planner who values
entrants' access to $w_m$ and the cross-jurisdiction price externality in \eqref{eq:price}. This
gap is the model's answer to Question 1, and it is exactly the inefficiency the Decentralization
Theorem's spillover condition rules in.
\end{proposition}

\section{Reaction functions: the empirical signature of the spillover}\label{sec:react}

Because $p_{s(j)}$ depends on neighbors' regulation, \eqref{eq:owner} is a game, with a reaction
function $R_j(\bar z_{-j},\tau_{-j})$ on each margin.

\begin{proposition}[Regulation is strategically substitutable within a submarket]\label{prop:react}
For co-submarket jurisdictions, if a neighbor loosens (raises $\bar z_k$ or lowers $\tau_k$), the
submarket price falls, raising a homevoter-pivotal $j$'s incentive to restrict: $j$ tightens.
Neighbors' regulation are strategic substitutes; each free-rides on the others' supply, and the
Nash equilibrium has less regional construction than the internalized optimum.
\end{proposition}

\paragraph{Why this is the centerpiece.} The reaction function is not merely an empirical
regularity --- it is the \emph{direct test of the federalism breakdown}. Under the
Decentralization Theorem's efficient case (no spillovers), jurisdictions would not respond to one
another's regulation at all: each would serve its own residents independently and the
cross-jurisdiction slope would be zero. A \emph{nonzero} strategic-interaction slope is therefore
the spillover, made visible --- the empirical object whose sign and magnitude establish that the
Oates condition fails and quantify how strongly. This is why estimating $R_j$ is the core
empirical task, and why a multi-jurisdiction panel (not a single metro) is required:
single-jurisdiction data cannot, even in principle, reveal a reaction function.

\section{Preemption and its leakage}\label{sec:preempt}

State preemption (e.g.\ ministerial-approval mandates, by-right upzoning, the builder's remedy)
raises the marginal cost to $j$ of a tight \emph{envelope} --- forcing $\bar z_j$ up or
converting discretionary categories to ministerial ones --- while leaving the marginal cost of
the \emph{discretionary} tax $\tau_j$ unchanged.

\begin{proposition}[Leakage: substitution from envelope to discretion]\label{prop:leak}
With effective restriction $r_j=g(\bar z_j,\tau_j)$ and the two margins substitutes in producing
it, a preemption that raises $\bar z_j$ (the map loosens, as intended) induces a
homevoter-pivotal jurisdiction that still demands restriction to raise $\tau_j$ in compensation
(more items conditioned, more opposition entertained, more continuances). Effective restriction
falls by \emph{less} than the change in $\bar z_j$ alone implies, and \textbf{the conditional
approval rate need not move}. The \emph{leakage} is the share of intended loosening offset by the
discretionary margin: $0$ (preemption fully effective) to $1$ (fully clawed back). This is the
model's answer to Question 2.
\end{proposition}

\noindent Leakage is precisely what makes preemption a federalism question rather than a
mechanical policy: partial centralization fails not because centralization is wrong but because
it is \emph{incomplete}, leaving a decentralized margin through which the spillover reasserts
itself. Strategic substitutability (Prop.~\ref{prop:react}) compounds it: as preempted
jurisdictions claw back via $\tau$, co-submarket neighbors face a smaller price effect and adjust
in turn.

\section{Counterfactuals and the identification challenge}\label{sec:ident}

Both questions are counterfactual. Question 1 compares the observed decentralized equilibrium to
an internalized benchmark never observed; Question 2 asks how supply responds to a change in the
assignment of regulatory authority. Answering them requires recovering the structural primitives
--- the homevoter weights $\phi_j\lambda_j$, the disamenity $g(\cdot)$, the demand and
agglomeration parameters --- from \emph{equilibrium} behavior alone, since only the chosen
regulation is observed, not the objectives that generated it. This is the identification crux:
the reaction-function slope is consistent with several mechanisms (the spillover modeled here,
but also yardstick competition or common shocks; cf.\ the strategic-interaction literature
\citealp{brueckner1998testing,brueckner2003strategic}), and disentangling them is what licenses
the counterfactual. The empirical leverage comes from the \emph{preemption ramp itself}: the
staggered, rule-driven timing of state mandates (builder's-remedy exposure following
Housing-Element non-compliance) is a supply-side shock plausibly orthogonal to a jurisdiction's
contemporaneous demand shock, providing variation that moves $\bar z_j$ for reasons other than
local taste --- the wedge needed to separate the spillover from its observational rivals and to
identify the leakage of Prop.~\ref{prop:leak}.

\section{Mapping to data}\label{sec:data}

The by-right envelope $\bar z_j$ is read from zoning (and is largely available from existing
LLM-based measurement, so need not be rebuilt); the discretionary tax $\tau_j$ is built fresh
from planning-commission minutes, item by item, across Bay Area jurisdictions; the preemption
ramp is dated from HCD Housing-Element compliance records; prices and quantities come from
transaction data. The reaction-function estimation (Prop.~\ref{prop:react}) is a spatial-lag
regression of a jurisdiction's discretionary tax on co-submarket neighbors', instrumented using
the preemption variation of \S\ref{sec:ident}; the leakage estimation (Prop.~\ref{prop:leak}) is
a difference-in-differences around preemption exposure with the \emph{envelope} and the
\emph{discretionary tax} as separate outcomes --- predicting the map loosens while discretion
tightens and the approval rate is flat. A short panel straddling the 2017--2023 preemption ramp
suffices for both; the data infrastructure is documented separately.

\section{Scope and a noted extension}\label{sec:scope}

The questions above are \emph{static/short-panel} counterfactuals: the magnitude of the
fragmentation distortion and the leakage of preemption are identified off the cross-section of
jurisdictions and the preemption ramp, and do not require deep historical data. A natural but
\emph{separate} extension treats regulation as state-dependent --- a ``regulatory ratchet'' in
which restriction is path-dependent and preemption shocks have persistent or hysteretic effects,
propagating across co-submarket neighbors over time. That dynamic question would require a long
panel and confronts a harder identification problem (state dependence vs.\ unobserved
heterogeneity); it is deliberately held as an extension rather than a core result, since the two
headline questions do not depend on it.

\section{Conclusion}

Fragmented housing regulation is a textbook fiscal-federalism breakdown: a local public good
(low density) whose provision via regulation imposes regional spillovers, so decentralized
choice over-restricts. Two moves make the breakdown quantitative and policy-relevant ---
splitting regulation into a by-right envelope and a discretionary friction (so that preemption
\emph{leakage} is a coherent object), and embedding many jurisdictions in a game whose
reaction functions are the direct empirical signature of the spillover. The model answers two
questions a descriptive account cannot: how large the fragmentation distortion is, and how much
of state preemption localities claw back through the margin it leaves untouched.

\bibliographystyle{plainnat}
\bibliography{toy_model}

\end{document}